
\documentclass{article}
\usepackage{colortbl}
\usepackage{makecell}
\usepackage{multirow}
\usepackage{supertabular}

\begin{document}

\newcounter{utterance}

\centering \large Interaction Transcript for game `matchit\_ascii', experiment `similar\_grid\_1', episode 1 with gpt{-}5.4{-}azure.
\vspace{24pt}

{ \footnotesize  \setcounter{utterance}{1}
\setlength{\tabcolsep}{0pt}
\begin{supertabular}{c@{$\;$}|p{.15\linewidth}@{}p{.15\linewidth}p{.15\linewidth}p{.15\linewidth}p{.15\linewidth}p{.15\linewidth}}
    \# & $\;$A & \multicolumn{4}{c}{Game Master} & $\;\:$B\\
    \hline

    \theutterance \stepcounter{utterance}  
    & & \multicolumn{4}{p{0.6\linewidth}}{
        \cellcolor[rgb]{0.9,0.9,0.9}{
            \makecell[{{p{\linewidth}}}]{
                \texttt{\tiny{[P1$\langle$GM]}}
                \texttt{Dalyvaujate bendrame spėjimo žaidime.} \\
\texttt{Tikslas {-} išsiaiškinti, ar šis tinklelis ir kitas paveikslėlis, kurį matau tik aš, yra vienodi.} \\
\texttt{Pirmiausia apibūdinkite savo tinklelį. Tada aš pateiksiu savo aprašymą ir galėsime vienas kitam užduoti klausimų apie tinklelius, kad išsiaiškintume, ar jie vienodi. Dabar pradėkite savo trumpą tinklelio aprašymą žodžiais "APRAŠYMAS:", po kurių rašykite aprašymą. Nieko daugiau nepridėkite. Tai yra jūsų tinklelis:} \\
\texttt{X X X X X} \\
\texttt{X X X X ▢} \\
\texttt{X X X ▢ ▢} \\
\texttt{X X ▢ ▢ ▢} \\
\texttt{X ▢ ▢ ▢ ▢} \\
            }
        }
    }
    & \\ \\

    \theutterance \stepcounter{utterance}  
    & \multicolumn{4}{p{0.6\linewidth}}{
        \cellcolor[rgb]{0.9,0.9,0.9}{
            \makecell[{{p{\linewidth}}}]{
                \texttt{\tiny{[P1$\rangle$GM]}}
                \texttt{APRAŠYMAS: 5x5 tinklelis, kuriame X sudaro trikampį viršutiniame kairiajame kampe: pirmoje eilutėje 5 X, antroje 4 X ir 1 ▢, trečioje 3 X ir 2 ▢, ketvirtoje 2 X ir 3 ▢, penktoje 1 X ir 4 ▢.} \\
            }
        }
    }
    & & \\ \\

    \theutterance \stepcounter{utterance}  
    & & & \multicolumn{2}{p{0.3\linewidth}}{
        \cellcolor[rgb]{0.9,0.9,0.9}{
            \makecell[{{p{\linewidth}}}]{
                \texttt{\tiny{[GM$|$GM]}}
                \texttt{continue} \\
            }
        }
    }
    & & \\ \\

    \theutterance \stepcounter{utterance}  
    & & \multicolumn{4}{p{0.6\linewidth}}{
        \cellcolor[rgb]{0.9,0.9,0.9}{
            \makecell[{{p{\linewidth}}}]{
                \texttt{\tiny{[GM$\rangle$P2]}}
                \texttt{Dalyvaujate bendrame spėjimo žaidime.} \\
\texttt{Tikslas {-} išsiaiškinti, ar šis tinklelis ir kitas paveikslėlis, kurį matau tik aš, yra vienodi.} \\
\texttt{Pirmiausia apibūdinkite savo tinklelį. Tada aš pateiksiu savo aprašymą ir galėsime vienas kitam užduoti klausimų apie tinklelius, kad išsiaiškintume, ar jie vienodi. Dabar pradėkite savo trumpą tinklelio aprašymą žodžiais "APRAŠYMAS:", po kurių rašykite aprašymą. Nieko daugiau nepridėkite. Tai yra jūsų tinklelis:} \\
\texttt{X X X X X} \\
\texttt{▢ X X X X} \\
\texttt{▢ ▢ X X X} \\
\texttt{▢ ▢ ▢ X X} \\
\texttt{▢ ▢ ▢ ▢ X} \\
            }
        }
    }
    & \\ \\

    \theutterance \stepcounter{utterance}  
    & & & \multicolumn{4}{p{0.6\linewidth}}{
        \cellcolor[rgb]{0.9,0.9,0.9}{
            \makecell[{{p{\linewidth}}}]{
                \texttt{\tiny{[GM$\langle$P2]}}
                \texttt{APRAŠYMAS: 5x5 tinklelis su įstrižaine iš viršutinio kairiojo į apatinį dešinįjį kampą: visi langeliai ant įstrižainės ir virš jos pažymėti X, o visi žemiau jos – tušti.} \\
            }
        }
    }
     \\ \\

    \theutterance \stepcounter{utterance}  
    & & & \multicolumn{2}{p{0.3\linewidth}}{
        \cellcolor[rgb]{0.9,0.9,0.9}{
            \makecell[{{p{\linewidth}}}]{
                \texttt{\tiny{[GM$|$GM]}}
                \texttt{continue} \\
            }
        }
    }
    & & \\ \\

    \theutterance \stepcounter{utterance}  
    & & \multicolumn{4}{p{0.6\linewidth}}{
        \cellcolor[rgb]{0.9,0.9,0.9}{
            \makecell[{{p{\linewidth}}}]{
                \texttt{\tiny{[GM$\rangle$P2]}}
                \texttt{Tai mano APRAŠYMASAPRAŠYMAS: 5x5 tinklelis, kuriame X sudaro trikampį viršutiniame kairiajame kampe: pirmoje eilutėje 5 X, antroje 4 X ir 1 ▢, trečioje 3 X ir 2 ▢, ketvirtoje 2 X ir 3 ▢, penktoje 1 X ir 4 ▢.} \\
\texttt{Dabar užduokite klausimą, kad sužinotumėte naujus mano tinklelio aspektus, kurie gali skirtis nuo jūsų tinklelio. Pradėkite "KLAUSIMAS:" ir nieko daugiau nepridėkite.} \\
            }
        }
    }
    & \\ \\

    \theutterance \stepcounter{utterance}  
    & & & \multicolumn{4}{p{0.6\linewidth}}{
        \cellcolor[rgb]{0.9,0.9,0.9}{
            \makecell[{{p{\linewidth}}}]{
                \texttt{\tiny{[GM$\langle$P2]}}
                \texttt{KLAUSIMAS: Ar kiekvienoje eilutėje tušti langeliai yra dešinėje nuo X, ar kairėje?} \\
            }
        }
    }
     \\ \\

    \theutterance \stepcounter{utterance}  
    & & & \multicolumn{2}{p{0.3\linewidth}}{
        \cellcolor[rgb]{0.9,0.9,0.9}{
            \makecell[{{p{\linewidth}}}]{
                \texttt{\tiny{[GM$|$GM]}}
                \texttt{continue} \\
            }
        }
    }
    & & \\ \\

    \theutterance \stepcounter{utterance}  
    & & \multicolumn{4}{p{0.6\linewidth}}{
        \cellcolor[rgb]{0.9,0.9,0.9}{
            \makecell[{{p{\linewidth}}}]{
                \texttt{\tiny{[P1$\langle$GM]}}
                \texttt{Tai mano APRAŠYMASAPRAŠYMAS: 5x5 tinklelis su įstrižaine iš viršutinio kairiojo į apatinį dešinįjį kampą: visi langeliai ant įstrižainės ir virš jos pažymėti X, o visi žemiau jos – tušti.} \\
\texttt{KLAUSIMAS: Ar kiekvienoje eilutėje tušti langeliai yra dešinėje nuo X, ar kairėje?Pradėkite atsakymą „ATSAKYMAS:" ir nieko daugiau nepridėkite.} \\
            }
        }
    }
    & \\ \\

    \theutterance \stepcounter{utterance}  
    & \multicolumn{4}{p{0.6\linewidth}}{
        \cellcolor[rgb]{0.9,0.9,0.9}{
            \makecell[{{p{\linewidth}}}]{
                \texttt{\tiny{[P1$\rangle$GM]}}
                \texttt{ATSAKYMAS: Kiekvienoje eilutėje tušti langeliai yra dešinėje nuo X.} \\
            }
        }
    }
    & & \\ \\

    \theutterance \stepcounter{utterance}  
    & & & \multicolumn{2}{p{0.3\linewidth}}{
        \cellcolor[rgb]{0.9,0.9,0.9}{
            \makecell[{{p{\linewidth}}}]{
                \texttt{\tiny{[GM$|$GM]}}
                \texttt{continue} \\
            }
        }
    }
    & & \\ \\

    \theutterance \stepcounter{utterance}  
    & & \multicolumn{4}{p{0.6\linewidth}}{
        \cellcolor[rgb]{0.9,0.9,0.9}{
            \makecell[{{p{\linewidth}}}]{
                \texttt{\tiny{[P1$\langle$GM]}}
                \texttt{Dabar užduokite klausimą, kad sužinotumėte naujus mano tinklelio aspektus, kurie gali skirtis nuo jūsų tinklelio. Pradėkite "KLAUSIMAS:" ir nieko daugiau nepridėkite.} \\
            }
        }
    }
    & \\ \\

    \theutterance \stepcounter{utterance}  
    & \multicolumn{4}{p{0.6\linewidth}}{
        \cellcolor[rgb]{0.9,0.9,0.9}{
            \makecell[{{p{\linewidth}}}]{
                \texttt{\tiny{[P1$\rangle$GM]}}
                \texttt{KLAUSIMAS: Ar kiekvienoje eilutėje X sudaro vientisą atkarpą nuo kairiojo krašto?} \\
            }
        }
    }
    & & \\ \\

    \theutterance \stepcounter{utterance}  
    & & & \multicolumn{2}{p{0.3\linewidth}}{
        \cellcolor[rgb]{0.9,0.9,0.9}{
            \makecell[{{p{\linewidth}}}]{
                \texttt{\tiny{[GM$|$GM]}}
                \texttt{continue} \\
            }
        }
    }
    & & \\ \\

    \theutterance \stepcounter{utterance}  
    & & \multicolumn{4}{p{0.6\linewidth}}{
        \cellcolor[rgb]{0.9,0.9,0.9}{
            \makecell[{{p{\linewidth}}}]{
                \texttt{\tiny{[GM$\rangle$P2]}}
                \texttt{ATSAKYMAS: Kiekvienoje eilutėje tušti langeliai yra dešinėje nuo X.} \\
\texttt{KLAUSIMAS: Ar kiekvienoje eilutėje X sudaro vientisą atkarpą nuo kairiojo krašto?Pradėkite atsakymą „ATSAKYMAS:" ir nieko daugiau nepridėkite.} \\
            }
        }
    }
    & \\ \\

    \theutterance \stepcounter{utterance}  
    & & & \multicolumn{4}{p{0.6\linewidth}}{
        \cellcolor[rgb]{0.9,0.9,0.9}{
            \makecell[{{p{\linewidth}}}]{
                \texttt{\tiny{[GM$\langle$P2]}}
                \texttt{ATSAKYMAS: Ne, tik pirmoje eilutėje X sudaro vientisą atkarpą nuo kairiojo krašto.} \\
            }
        }
    }
     \\ \\

    \theutterance \stepcounter{utterance}  
    & & & \multicolumn{2}{p{0.3\linewidth}}{
        \cellcolor[rgb]{0.9,0.9,0.9}{
            \makecell[{{p{\linewidth}}}]{
                \texttt{\tiny{[GM$|$GM]}}
                \texttt{continue} \\
            }
        }
    }
    & & \\ \\

    \theutterance \stepcounter{utterance}  
    & & \multicolumn{4}{p{0.6\linewidth}}{
        \cellcolor[rgb]{0.9,0.9,0.9}{
            \makecell[{{p{\linewidth}}}]{
                \texttt{\tiny{[GM$\rangle$P2]}}
                \texttt{Dabar užduokite klausimą, kad sužinotumėte naujus mano tinklelio aspektus, kurie gali skirtis nuo jūsų tinklelio. Pradėkite "KLAUSIMAS:" ir nieko daugiau nepridėkite.} \\
            }
        }
    }
    & \\ \\

    \theutterance \stepcounter{utterance}  
    & & & \multicolumn{4}{p{0.6\linewidth}}{
        \cellcolor[rgb]{0.9,0.9,0.9}{
            \makecell[{{p{\linewidth}}}]{
                \texttt{\tiny{[GM$\langle$P2]}}
                \texttt{KLAUSIMAS: Ar X jūsų tinklelyje yra ant pagrindinės įstrižainės nuo viršutinio kairiojo iki apatinio dešiniojo kampo?} \\
            }
        }
    }
     \\ \\

    \theutterance \stepcounter{utterance}  
    & & & \multicolumn{2}{p{0.3\linewidth}}{
        \cellcolor[rgb]{0.9,0.9,0.9}{
            \makecell[{{p{\linewidth}}}]{
                \texttt{\tiny{[GM$|$GM]}}
                \texttt{continue} \\
            }
        }
    }
    & & \\ \\

    \theutterance \stepcounter{utterance}  
    & & \multicolumn{4}{p{0.6\linewidth}}{
        \cellcolor[rgb]{0.9,0.9,0.9}{
            \makecell[{{p{\linewidth}}}]{
                \texttt{\tiny{[P1$\langle$GM]}}
                \texttt{ATSAKYMAS: Ne, tik pirmoje eilutėje X sudaro vientisą atkarpą nuo kairiojo krašto.} \\
\texttt{KLAUSIMAS: Ar X jūsų tinklelyje yra ant pagrindinės įstrižainės nuo viršutinio kairiojo iki apatinio dešiniojo kampo?Pradėkite atsakymą „ATSAKYMAS:" ir nieko daugiau nepridėkite.} \\
            }
        }
    }
    & \\ \\

    \theutterance \stepcounter{utterance}  
    & \multicolumn{4}{p{0.6\linewidth}}{
        \cellcolor[rgb]{0.9,0.9,0.9}{
            \makecell[{{p{\linewidth}}}]{
                \texttt{\tiny{[P1$\rangle$GM]}}
                \texttt{ATSAKYMAS: Taip, bet tik pirmuose trijuose pagrindinės įstrižainės langeliuose.} \\
            }
        }
    }
    & & \\ \\

    \theutterance \stepcounter{utterance}  
    & & & \multicolumn{2}{p{0.3\linewidth}}{
        \cellcolor[rgb]{0.9,0.9,0.9}{
            \makecell[{{p{\linewidth}}}]{
                \texttt{\tiny{[GM$|$GM]}}
                \texttt{continue} \\
            }
        }
    }
    & & \\ \\

    \theutterance \stepcounter{utterance}  
    & & \multicolumn{4}{p{0.6\linewidth}}{
        \cellcolor[rgb]{0.9,0.9,0.9}{
            \makecell[{{p{\linewidth}}}]{
                \texttt{\tiny{[P1$\langle$GM]}}
                \texttt{Dabar užduokite klausimą, kad sužinotumėte naujus mano tinklelio aspektus, kurie gali skirtis nuo jūsų tinklelio. Pradėkite "KLAUSIMAS:" ir nieko daugiau nepridėkite.} \\
            }
        }
    }
    & \\ \\

    \theutterance \stepcounter{utterance}  
    & \multicolumn{4}{p{0.6\linewidth}}{
        \cellcolor[rgb]{0.9,0.9,0.9}{
            \makecell[{{p{\linewidth}}}]{
                \texttt{\tiny{[P1$\rangle$GM]}}
                \texttt{KLAUSIMAS: Ar kiekvienoje eilutėje X sudaro vientisą atkarpą iki dešiniojo krašto?} \\
            }
        }
    }
    & & \\ \\

    \theutterance \stepcounter{utterance}  
    & & & \multicolumn{2}{p{0.3\linewidth}}{
        \cellcolor[rgb]{0.9,0.9,0.9}{
            \makecell[{{p{\linewidth}}}]{
                \texttt{\tiny{[GM$|$GM]}}
                \texttt{continue} \\
            }
        }
    }
    & & \\ \\

    \theutterance \stepcounter{utterance}  
    & & \multicolumn{4}{p{0.6\linewidth}}{
        \cellcolor[rgb]{0.9,0.9,0.9}{
            \makecell[{{p{\linewidth}}}]{
                \texttt{\tiny{[GM$\rangle$P2]}}
                \texttt{ATSAKYMAS: Taip, bet tik pirmuose trijuose pagrindinės įstrižainės langeliuose.} \\
\texttt{KLAUSIMAS: Ar kiekvienoje eilutėje X sudaro vientisą atkarpą iki dešiniojo krašto?Pradėkite atsakymą „ATSAKYMAS:" ir nieko daugiau nepridėkite.} \\
            }
        }
    }
    & \\ \\

    \theutterance \stepcounter{utterance}  
    & & & \multicolumn{4}{p{0.6\linewidth}}{
        \cellcolor[rgb]{0.9,0.9,0.9}{
            \makecell[{{p{\linewidth}}}]{
                \texttt{\tiny{[GM$\langle$P2]}}
                \texttt{ATSAKYMAS: Taip, kiekvienoje eilutėje X sudaro vientisą atkarpą iki dešiniojo krašto.} \\
            }
        }
    }
     \\ \\

    \theutterance \stepcounter{utterance}  
    & & & \multicolumn{2}{p{0.3\linewidth}}{
        \cellcolor[rgb]{0.9,0.9,0.9}{
            \makecell[{{p{\linewidth}}}]{
                \texttt{\tiny{[GM$|$GM]}}
                \texttt{continue} \\
            }
        }
    }
    & & \\ \\

    \theutterance \stepcounter{utterance}  
    & & \multicolumn{4}{p{0.6\linewidth}}{
        \cellcolor[rgb]{0.9,0.9,0.9}{
            \makecell[{{p{\linewidth}}}]{
                \texttt{\tiny{[GM$\rangle$P2]}}
                \texttt{Dabar užduokite klausimą, kad sužinotumėte naujus mano tinklelio aspektus, kurie gali skirtis nuo jūsų tinklelio. Pradėkite "KLAUSIMAS:" ir nieko daugiau nepridėkite.} \\
            }
        }
    }
    & \\ \\

    \theutterance \stepcounter{utterance}  
    & & & \multicolumn{4}{p{0.6\linewidth}}{
        \cellcolor[rgb]{0.9,0.9,0.9}{
            \makecell[{{p{\linewidth}}}]{
                \texttt{\tiny{[GM$\langle$P2]}}
                \texttt{KLAUSIMAS: Ar apatiniame dešiniajame kampe yra X?} \\
            }
        }
    }
     \\ \\

    \theutterance \stepcounter{utterance}  
    & & & \multicolumn{2}{p{0.3\linewidth}}{
        \cellcolor[rgb]{0.9,0.9,0.9}{
            \makecell[{{p{\linewidth}}}]{
                \texttt{\tiny{[GM$|$GM]}}
                \texttt{continue} \\
            }
        }
    }
    & & \\ \\

    \theutterance \stepcounter{utterance}  
    & & \multicolumn{4}{p{0.6\linewidth}}{
        \cellcolor[rgb]{0.9,0.9,0.9}{
            \makecell[{{p{\linewidth}}}]{
                \texttt{\tiny{[P1$\langle$GM]}}
                \texttt{ATSAKYMAS: Taip, kiekvienoje eilutėje X sudaro vientisą atkarpą iki dešiniojo krašto.} \\
\texttt{KLAUSIMAS: Ar apatiniame dešiniajame kampe yra X?Pradėkite atsakymą „ATSAKYMAS:" ir nieko daugiau nepridėkite.} \\
            }
        }
    }
    & \\ \\

    \theutterance \stepcounter{utterance}  
    & \multicolumn{4}{p{0.6\linewidth}}{
        \cellcolor[rgb]{0.9,0.9,0.9}{
            \makecell[{{p{\linewidth}}}]{
                \texttt{\tiny{[P1$\rangle$GM]}}
                \texttt{ATSAKYMAS: Ne.} \\
            }
        }
    }
    & & \\ \\

    \theutterance \stepcounter{utterance}  
    & & & \multicolumn{2}{p{0.3\linewidth}}{
        \cellcolor[rgb]{0.9,0.9,0.9}{
            \makecell[{{p{\linewidth}}}]{
                \texttt{\tiny{[GM$|$GM]}}
                \texttt{continue} \\
            }
        }
    }
    & & \\ \\

    \theutterance \stepcounter{utterance}  
    & & \multicolumn{4}{p{0.6\linewidth}}{
        \cellcolor[rgb]{0.9,0.9,0.9}{
            \makecell[{{p{\linewidth}}}]{
                \texttt{\tiny{[P1$\langle$GM]}}
                \texttt{Dabar užduokite klausimą, kad sužinotumėte naujus mano tinklelio aspektus, kurie gali skirtis nuo jūsų tinklelio. Pradėkite "KLAUSIMAS:" ir nieko daugiau nepridėkite.} \\
            }
        }
    }
    & \\ \\

    \theutterance \stepcounter{utterance}  
    & \multicolumn{4}{p{0.6\linewidth}}{
        \cellcolor[rgb]{0.9,0.9,0.9}{
            \makecell[{{p{\linewidth}}}]{
                \texttt{\tiny{[P1$\rangle$GM]}}
                \texttt{KLAUSIMAS: Ar apatiniame kairiajame kampe yra X?} \\
            }
        }
    }
    & & \\ \\

    \theutterance \stepcounter{utterance}  
    & & & \multicolumn{2}{p{0.3\linewidth}}{
        \cellcolor[rgb]{0.9,0.9,0.9}{
            \makecell[{{p{\linewidth}}}]{
                \texttt{\tiny{[GM$|$GM]}}
                \texttt{continue} \\
            }
        }
    }
    & & \\ \\

    \theutterance \stepcounter{utterance}  
    & & \multicolumn{4}{p{0.6\linewidth}}{
        \cellcolor[rgb]{0.9,0.9,0.9}{
            \makecell[{{p{\linewidth}}}]{
                \texttt{\tiny{[GM$\rangle$P2]}}
                \texttt{ATSAKYMAS: Ne.} \\
\texttt{KLAUSIMAS: Ar apatiniame kairiajame kampe yra X?Pradėkite atsakymą „ATSAKYMAS:" ir nieko daugiau nepridėkite.} \\
            }
        }
    }
    & \\ \\

    \theutterance \stepcounter{utterance}  
    & & & \multicolumn{4}{p{0.6\linewidth}}{
        \cellcolor[rgb]{0.9,0.9,0.9}{
            \makecell[{{p{\linewidth}}}]{
                \texttt{\tiny{[GM$\langle$P2]}}
                \texttt{ATSAKYMAS: Ne.} \\
            }
        }
    }
     \\ \\

    \theutterance \stepcounter{utterance}  
    & & & \multicolumn{2}{p{0.3\linewidth}}{
        \cellcolor[rgb]{0.9,0.9,0.9}{
            \makecell[{{p{\linewidth}}}]{
                \texttt{\tiny{[GM$|$GM]}}
                \texttt{continue} \\
            }
        }
    }
    & & \\ \\

    \theutterance \stepcounter{utterance}  
    & & \multicolumn{4}{p{0.6\linewidth}}{
        \cellcolor[rgb]{0.9,0.9,0.9}{
            \makecell[{{p{\linewidth}}}]{
                \texttt{\tiny{[GM$\rangle$P2]}}
                \texttt{Dabar priimkite sprendimą. Kaip manote: ar jūsų tinklelis ir kitas aprašytas tinklelis yra tokie patys? Parašykite "SPRENDIMAS: tas pats tinklelis", jei manote, kad tai tas pats tinklelis, arba "SPRENDIMAS: skirtingi tinkleliai", jei manote, kad tai skirtingi tinkleliai. Nieko daugiau nepridėkite.} \\
            }
        }
    }
    & \\ \\

    \theutterance \stepcounter{utterance}  
    & & & \multicolumn{4}{p{0.6\linewidth}}{
        \cellcolor[rgb]{0.9,0.9,0.9}{
            \makecell[{{p{\linewidth}}}]{
                \texttt{\tiny{[GM$\langle$P2]}}
                \texttt{SPRENDIMAS: skirtingi tinkleliai} \\
            }
        }
    }
     \\ \\

    \theutterance \stepcounter{utterance}  
    & & & \multicolumn{2}{p{0.3\linewidth}}{
        \cellcolor[rgb]{0.9,0.9,0.9}{
            \makecell[{{p{\linewidth}}}]{
                \texttt{\tiny{[GM$|$GM]}}
                \texttt{continue} \\
            }
        }
    }
    & & \\ \\

    \theutterance \stepcounter{utterance}  
    & & & \multicolumn{2}{p{0.3\linewidth}}{
        \cellcolor[rgb]{0.9,0.9,0.9}{
            \makecell[{{p{\linewidth}}}]{
                \texttt{\tiny{[GM$|$GM]}}
                \texttt{success} \\
            }
        }
    }
    & & \\ \\

    \theutterance \stepcounter{utterance}  
    & & \multicolumn{4}{p{0.6\linewidth}}{
        \cellcolor[rgb]{0.9,0.9,0.9}{
            \makecell[{{p{\linewidth}}}]{
                \texttt{\tiny{[P1$\langle$GM]}}
                \texttt{ATSAKYMAS: Ne.} \\
\texttt{Dabar priimkite sprendimą. Kaip manote: ar jūsų tinklelis ir kitas aprašytas tinklelis yra tokie patys? Parašykite "SPRENDIMAS: tas pats tinklelis", jei manote, kad tai tas pats tinklelis, arba "SPRENDIMAS: skirtingi tinkleliai", jei manote, kad tai skirtingi tinkleliai. Nieko daugiau nepridėkite.} \\
            }
        }
    }
    & \\ \\

    \theutterance \stepcounter{utterance}  
    & \multicolumn{4}{p{0.6\linewidth}}{
        \cellcolor[rgb]{0.9,0.9,0.9}{
            \makecell[{{p{\linewidth}}}]{
                \texttt{\tiny{[P1$\rangle$GM]}}
                \texttt{SPRENDIMAS: skirtingi tinkleliai} \\
            }
        }
    }
    & & \\ \\

    \theutterance \stepcounter{utterance}  
    & & & \multicolumn{2}{p{0.3\linewidth}}{
        \cellcolor[rgb]{0.9,0.9,0.9}{
            \makecell[{{p{\linewidth}}}]{
                \texttt{\tiny{[GM$|$GM]}}
                \texttt{continue} \\
            }
        }
    }
    & & \\ \\

    \theutterance \stepcounter{utterance}  
    & & & \multicolumn{2}{p{0.3\linewidth}}{
        \cellcolor[rgb]{0.9,0.9,0.9}{
            \makecell[{{p{\linewidth}}}]{
                \texttt{\tiny{[GM$|$GM]}}
                \texttt{success} \\
            }
        }
    }
    & & \\ \\

\end{supertabular}
}

\end{document}
